\subsection{Описание предметной области}
\label{sec:analysis:domain}

Рынок автомобильных запасных частей и компонентов образует динамичный сектор розничной электронной коммерции. В отличие от типовых потребительских товаров, номенклатура автозапчастей насчитывает сотни тысяч наименований и охватывает широкий спектр категорий: детали двигателя, элементы подвески, тормозной системы, трансмиссии, электрооборудование и расходные материалы. Наряду с оригинальными компонентами на рынке обращаются лицензионные аналоги сторонних брендов с различными характеристиками и ценами [1].

В бизнес-процессах интернет-магазина автозапчастей участвуют две основные группы пользователей: покупатели и администраторы. Покупателям (автовладельцам и мастерам сервисных центров) необходим оперативный поиск деталей по марке и модели автомобиля, бренду и категории, добавление позиций в корзину и быстрое оформление заказа. Администраторы и складские менеджеры выполняют ведение номенклатуры, контроль складских запасов, управление статусами заказов, импорт и экспорт каталогов, а в завершение рабочих циклов формируют аналитическую отчетность.

Ведение торговли без средств автоматизации приводит к существенным издержкам. Ручной учет сопровождается длительным поиском совместимых деталей среди обширного ассортимента и риском продажи отсутствующего товара при одновременном оформлении одного физического экземпляра несколькими покупателями. Изменение базовых цен номенклатуры задним числом искажает финансовую отчетность из-за отсутствия фиксации цен сделки в момент заказа. Ручное обновление прайс-листов от внешних поставщиков отличается высокой трудоемкостью, а руководству не хватает оперативной аналитики по выручке, стадиям заказов и наиболее востребованным позициям.

Формализация предметной области в терминах реляционных данных определяет восемь взаимосвязанных сущностей. Сущность «Роль» разграничивает уровни полномочий в системе, а сущность «Пользователь» хранит учетные записи клиентов и администраторов. Сведения о запчастях сосредоточены в сущности «Товар», содержащей наименование, бренд, модель автомобиля, категорию, цену и складской остаток. Выбранные покупателем позиции временно объединяются в сущности «Корзина». Согласованная сделка фиксируется сущностью «Заказ», строки которого образуют сущность «Позиция заказа» с фиксацией контрактной стоимости и количества деталей. Финансовая сторона сделки отражается сущностью «Транзакция», а адресные реквизиты и логистический статус фиксируются сущностью «Данные доставки».

Главный процесс системы заключается в транзакционном переводе выбранных товаров из корзины в оформленный заказ. В рамках единой транзакции СУБД проверяется наличие необходимого количества деталей на складе, фиксируются закупочные цены позиций, списываются складские запасы и регистрируются реквизиты доставки и оплаты. При отмене заказа система автоматически возвращает зарезервированные товары на склад.
